% from ega4-11-fr.tex
\oldpage[IV-3]{143}

soit bijective, puisque $C$ est une $B$-algèbre de type fini. Mais alors
l'application $B_\lambda\to B$ est bijective par fidèle platitude, donc
$B=B_\lambda$ est une $A$-algèbre de type fini.

Supposons maintenant que $g\circ f$ soit de présentation finie, $C$ étant donc
une $A$-algèbre de présentation finie. D'après la première partie du
raisonnement, on a
\[
B=A[T_1,\ldots,T_n]/\widetilde{\mathfrak J}
\]
pour un idéal $\widetilde{\mathfrak J}$~; soit
$(\widetilde{\mathfrak J}_\lambda)$ la famille filtrante des idéaux de type
fini de $A[T_1,\ldots,T_n]$ contenus dans $\widetilde{\mathfrak J}$, de sorte
que $\widetilde{\mathfrak J}=\varinjlim\widetilde{\mathfrak J}_\lambda$ et
$B=\varinjlim B_\lambda$, avec
$B_\lambda=A[T_1,\ldots,T_n]/\widetilde{\mathfrak J}_\lambda$. Appliquant
comme ci-dessus (8.8.2), (8.10.5) et (11.2.6), on a
$C=C_\alpha\otimes_{B_\alpha}B$, où $C_\alpha$ est une
$B_\alpha$-algèbre de présentation finie et un $B_\alpha$-module fidèlement
plat~; on posera encore $C_\lambda=C_\alpha\otimes_{B_\alpha}B_\lambda$ pour
$\lambda\geq\alpha$, de sorte que
\[
C_\lambda=C_\alpha/\big(C_\alpha\otimes_{B_\alpha}
(\widetilde{\mathfrak J}_\lambda/\widetilde{\mathfrak J}_\alpha)\big)
\]
par platitude, et de même
\[
C=C_\alpha/\big(C_\alpha\otimes_{B_\alpha}
(\widetilde{\mathfrak J}/\widetilde{\mathfrak J}_\alpha)\big).
\]
Comme par hypothèse $C$ est une $A$-algèbre de présentation finie ainsi que
$C_\alpha$, $C_\alpha\otimes_{B_\alpha}
(\widetilde{\mathfrak J}/\widetilde{\mathfrak J}_\alpha)$ est un idéal de
type fini de $C_\alpha$ (1.4.4), et les
$C_\alpha\otimes_{B_\alpha}
(\widetilde{\mathfrak J}_\lambda/\widetilde{\mathfrak J}_\alpha)$
s'identifient à des idéaux de type fini de $C_\alpha$ dont
$C_\alpha\otimes_{B_\alpha}
(\widetilde{\mathfrak J}/\widetilde{\mathfrak J}_\alpha)$ est la réunion. Il
existe par suite un $\lambda\geq\alpha$ tel que
\[
C_\alpha\otimes_{B_\alpha}
(\widetilde{\mathfrak J}/\widetilde{\mathfrak J}_\alpha)
=C_\alpha\otimes_{B_\alpha}
(\widetilde{\mathfrak J}_\lambda/\widetilde{\mathfrak J}_\alpha),
\]
d'où $\widetilde{\mathfrak J}=\widetilde{\mathfrak J}_\lambda$ par fidèle
platitude, ce qui prouve que $B$ est une $A$-algèbre de présentation finie.
\end{proof}

\subsection{Descente de la platitude par des morphismes quelconques~: cas d'un préschéma de base artinien}
\label{subsection:IV.11.4-fr}

\begin{theorem}[11.4.1]
\label{IV.11.4.1-fr}
Soient $A$ un anneau, $\mathfrak J$ un idéal nilpotent de $A$,
$u_\lambda:A\to B_\lambda$ $(\lambda\in L)$ une famille d'homomorphismes
d'anneaux telle que l'intersection des noyaux des $u_\lambda$ soit réduite à
$0$. Soit $M$ un $A$-module tel que pour tout $\lambda\in L$,
$M\otimes_A B_\lambda$ soit un $B_\lambda$-module libre et que
$M\otimes_A(A/\mathfrak J)$ soit un $(A/\mathfrak J)$-module libre. Alors $M$
est un $A$-module libre. Si de plus l'ensemble d'indices $L$ est fini, on peut
remplacer partout «~libre~» par «~plat~» dans l'énoncé précédent.
\end{theorem}

\begin{proof}
Dans les deux cas il suffira de prouver que $M$ est un $A$-module \emph{plat}~:
en effet, lorsque $M\otimes_A(A/\mathfrak J)$ est un
$(A/\mathfrak J)$-module libre, il en résultera que $M$ est un $A$-module libre
en vertu de $(\mathbf 0_{\mathrm{III}},10.1.2)$.

Nous utiliserons le lemme suivant, qui généralise
$(\mathbf 0_{\mathrm{III}},10.2.1)$~:

\begin{lemma}[11.4.1.1]
\label{IV.11.4.1.1-fr}
Soit $A$ un anneau muni d'une filtration finie
$(\mathfrak J_n)_{0\leq n\leq N+1}$, avec $\mathfrak J_0=A$,
$\mathfrak J_{N+1}=0$. Soit $M$ un $A$-module muni de la filtration
$(\mathfrak J_nM)_{0\leq n\leq N+1}$, et désignons par
$\operatorname{gr}(A)$ et $\operatorname{gr}(M)$ l'anneau et le module gradués
correspondants. Supposons que $M\otimes_A(A/\mathfrak J_1)$ soit un
$(A/\mathfrak J_1)$-module plat et que l'homomorphisme canonique
\begin{equation}
\operatorname{gr}_0(M)\otimes_{\operatorname{gr}_0(A)}\operatorname{gr}(A)
\longrightarrow\operatorname{gr}(M)
\label{IV.11.4.1.2-fr}
\tag{11.4.1.2}
\end{equation}
soit injectif. Alors $M$ est un $A$-module plat.
\end{lemma}

L'homomorphisme canonique (11.4.1.2) se définit de la même façon que celui de
$(\mathbf 0_{\mathrm{III}},10.1.1)$ comme étant en degré $n$ l'homomorphisme
\[
(M/\mathfrak J_1M)\otimes(\mathfrak J_n/\mathfrak J_{n+1})
=(M\otimes\mathfrak J_n)/
(\operatorname{Im}(M\otimes\mathfrak J_{n+1})
+\operatorname{Im}(\mathfrak J_1M\otimes\mathfrak J_n))
\longrightarrow\mathfrak J_nM/\mathfrak J_{n+1}M
\]
provenant de l'homomorphisme canonique $M\otimes\mathfrak J_n\to\mathfrak J_nM$
par passage aux quotients. Le lemme se démontre par récurrence sur $N$,
puisqu'il n'y a rien à démontrer pour $N=0$.

